% ===================================================================
%  圖表第 13 號 — 酒店。--- 餐廳。咖啡室。
%
%  Traduction, faite en 2026, en cantonais ecrit d'aujourd'hui, en
%  caracteres traditionnels, sur l'ido de texto/io. Regles de la
%  colonne : voir 10-toubiu-01.tex. Une seule numerotation. Deux
%  notes d'auteur, toutes deux en gras-italique dans le fac-simile.
%
%  L'HOTEL EST LE TABLEAU LE PLUS ANGLAIS DE LA COLONNE, ET IL PORTE
%  UN CARACTERE QUI N'EXISTE QU'A HONG KONG :
%
%    𨋢 (9)      l'ascenseur  Pekin : 电梯
%
%  Ce caractere-la n'est pas un mot chinois habille : il a ete FORGE
%  pour l'anglais « lift », en collant 車 a 卑, et il ne s'ecrit
%  nulle part ailleurs. Il n'est pas au repertoire de base d'Unicode,
%  il est dans les plans complementaires — un caractere que le nord
%  n'a jamais eu a coder.
%
%  Le reste du sejour est du meme bois :
%
%    貼士 (59)  le pourboire  Pekin : 小费    — « tips »
%    啤牌 (82)  les cartes    Pekin : 纸牌
%    喼 (27)    la valise     Pekin : 箱子
%    骨牌 (80)  les dominos   Pekin : 多米诺
%    收銀 (15)  la caisse     Pekin : 收款处
%    大堂 (68)  le hall
%    地庫 (48)  le sous-sol   Pekin : 地下室
%    煙仔 (8)   la cigarette  Pekin : 香烟
%    生果       les fruits    Pekin : 水果
%    洗衫 (42)  la lessive    Pekin : 洗衣
%    擔挑 (40)  la palanche
%
%  ET LES TROIS REPAS SE SEPARENT ICI COMME AU QUEBEC, MAIS AUTREMENT.
%  L'ido distingue « dejuneto » et « dejuno » dans le meme tableau —
%  c'est ce qui avait fait la preuve pour fr-CA. Le cantonais rend le
%  premier par 早餐 et le second par 晏 : 食晏, prendre le repas de
%  midi, est un verbe a lui seul, que le nord ne dit pas.
%
%  LES ETAGES SUIVENT HONG KONG, ET HONG KONG SUIT LE FRANCAIS :
%  地下 le rez-de-chaussee, 一樓 le premier, 二樓 (1) le deuxieme.
%  C'est le meme decompte qu'au tableau 6, et c'est justement celui
%  ou l'usage parle du Quebec s'ecartait.
%
%  QUATRE BLOCS ONT DEMANDE L'INVERSION OU L'APPOSITION : la chambre
%  et son etage (t13-01-2), le cocher de l'omnibus et la malle du
%  voyageur (t13-05-1), les nuages de l'horizon et l'enseigne du hall
%  (t13-10-1).
% ===================================================================

%%K t13-noto-1 noto
\VUnotes{143.2mm}{%
\textit{\VUgras{(*) 房入面嘅細節，睇圖表第 6 號。}}%
}

%%K t13-apar-1 apar
\VUsaut{3.0mm}
\VUtitre{15.0pt}{60}{圖表第 13 號}
\VUsaut{1.6mm}
\VUpk{10.2pt}{40}{酒店。--- 餐廳。}
\VUsaut{2.0mm}
\VUpk{10.2pt}{40}{咖啡室。}
\VUsaut{3.4mm}

%%K t13-01-1 p
1. --- 尋晚到咗魯瓦揚，我就住咗入「\textit{大洋酒店}」——係個朋友介紹我嘅。

%%K t13-01-2 p
佢畀咗一間靚\VUgras{房} \textsuperscript{(2)} 我，喺\VUgras{二樓} \textsuperscript{(1)}；
我瞓得好好 (*)。

%%K t13-02-1 p
2. --- 今朝行出間房——\VUgras{房務員} \textsuperscript{(3)} 已經送咗早餐入嚟——
我入咗個\VUgras{吸煙室} \textsuperscript{(5)}，嗰度同時做\VUgras{閱報室}
\textsuperscript{(4)}；我一路睇\VUgras{報紙} \textsuperscript{(6)}，一路食
\VUgras{煙仔} \textsuperscript{(8)}。睇完，我搭\VUgras{𨋢} \textsuperscript{(9)} 落樓，出去城入面行下。

%%K t13-03-1 p
3. --- 我唔記得咗問酒店個\VUgras{賬房} \textsuperscript{(12)} 幾點開\VUgras{圍檯}
\textsuperscript{(11)} 飯，所以早早就返咗嚟，順便去酒店個\VUgras{浴室}
\textsuperscript{(10)} 沖個涼。之後我又去\VUgras{收銀處} \textsuperscript{(15)} 打電話
畀一個朋友。點知部\VUgras{電話} \textsuperscript{(13)} 有人用緊；個\VUgras{女收銀}
\textsuperscript{(14)} 見我為難——佢當時正入緊張\VUgras{單} \textsuperscript{(17)}，入落
本\VUgras{旅客登記簿} \textsuperscript{(16)} 度——就叫個\VUgras{看門} \textsuperscript{(18)} 派
個\VUgras{後生仔} \textsuperscript{(19)} \VUgras{踩單車} \textsuperscript{(20)} 幫我走一趟。

%%K t13-04-1 p
4. --- 我多謝咗佢，就行出去打賞個\VUgras{搬運佬} \textsuperscript{(24)}（做粗重嘅
嗰個）——

%%K t13-04-1 p suite
佢用架\VUgras{手推車} \textsuperscript{(25)} 由火車站幫我運咗個\VUgras{帽盒}
\textsuperscript{(26)}、\VUgras{我個喼} \textsuperscript{(27)} 同\VUgras{我個旅行袋}
\textsuperscript{(28)} 過嚟。個\VUgras{郵差} \textsuperscript{(21)} 啱啱到，交咗封
\VUgras{信} \textsuperscript{(22)} 畀我。

%%K t13-05-1 p
5. --- 我喺門檻度仲企多陣，就喺個\VUgras{郵筒} \textsuperscript{(23)} 隔籬，睇下
街上面點行點走。有個\VUgras{車伕} \textsuperscript{(30)}——\VUgras{公共馬車}
\textsuperscript{(29)} 嗰個——正將個\VUgras{大喼} \textsuperscript{(31)} 搬上車，個喼係個
\VUgras{旅客} \textsuperscript{(32)} 嘅，\VUgras{酒店老闆} \textsuperscript{(33)} 就喺度送客。
一個後生\VUgras{電報生} \textsuperscript{(35)} 慢條斯理噉送緊份\VUgras{電報}
\textsuperscript{(34)} 畀酒店一位住客；有個\VUgras{遊客} \textsuperscript{(36)} 孭住部
\VUgras{相機} \textsuperscript{(38)}，喺度揸住本\VUgras{旅遊指南} \textsuperscript{(37)} 睇。

%%K t13-tit-1 sub
\VUsaut{4.0mm}
\VUpk{10.2pt}{40}{食晏時分。}
\VUsaut{3.0mm}

%%K t13-06-1 p
6. --- 鐵閘前面，有個唔憂唔慮嘅\VUgras{跑腿} \textsuperscript{(39)} 喺自己條
\VUgras{擔挑} \textsuperscript{(40)} 同個\VUgras{擦鞋箱} \textsuperscript{(41)} 中間瞌眼瞓；
一個後生\VUgras{洗衫妹} \textsuperscript{(42)} 孭住\VUgras{籮} \textsuperscript{(43)} 衫，見到
企喺門邊嗰個\VUgras{餐廳領班} \textsuperscript{(44)} 一副正經樣，就笑咗出嚟。

%%K t13-07-1 p
7. --- 我見到個\VUgras{大廚} \textsuperscript{(45)} 行出佢個\VUgras{廚房}
\textsuperscript{(47)}——廚房喺\VUgras{地庫} \textsuperscript{(48)}——出嚟叫個
\VUgras{洗碗嘅} \textsuperscript{(46)} 去辦件事；我就記起就快夠鐘食晏，於是行入餐廳。
穿過個天井嗰陣，有個\VUgras{司機} \textsuperscript{(50)} 泊緊佢架\VUgras{汽車}
\textsuperscript{(49)}；仲有個\VUgras{機工} \textsuperscript{(52)} 照住個\VUgras{踩單車嘅}
\textsuperscript{(53)} 講嘅去整架啱啱撞咗嘅\VUgras{汽油三輪車} \textsuperscript{(51)}。

%%K t13-08-1 p
8. --- 有個後生\VUgras{侍仔} \textsuperscript{(54)} 攞住幾杯\VUgras{啤酒}
\textsuperscript{(55)}，我就問佢圍檯飯係咪擺喺露台底下；

%%K t13-noto-2 noto
\VUnotes{143.2mm}{%
\textit{\VUgras{(*) 有咗字典入面嗰張菜式表，老師想點改下面呢張菜牌都好易。}}%
}

%%K t13-08-1 p suite
佢話係，我就揀咗其中一張檯坐低，叫咗一餐好嘢食。

%%K t13-tit-2 sub
\VUsaut{4.0mm}
\VUpk{10.2pt}{40}{一餐靚晏。}
\VUsaut{3.0mm}

%%K t13-09-1 p
9. --- 侍應攞咗張晏嘅\VUgras{菜牌} (*) 同酒牌畀我。我揀好之後，叫咗
\VUgras{蝦仔}配\VUgras{牛油}做\VUgras{頭盤，}前菜就叫咗\VUgras{卡昂式牛雜。}我冇食
\VUgras{魚，}要咗一份羊\VUgras{髀，}\VUgras{菜式}就叫咗忌廉\VUgras{菠菜。}跟住，
啲\VUgras{甜品}冇一款啱我心水，我就叫人送幾樣飯後嘢上嚟：\VUgras{芝士、}
\VUgras{生果}同\VUgras{餅乾。}我仲加咗一支好正嘅聖埃美隆\VUgras{紅酒。}
食到心滿意足，我畀咗個\VUgras{侍應} \textsuperscript{(57)} 封靚\VUgras{貼士}
\textsuperscript{(59)}，就離開張檯。

%%K t13-10-1 p
10. --- 見我打算走，個\VUgras{經理} \textsuperscript{(60)} 就話不如上露台睇下
\VUgras{大海} \textsuperscript{(62)} 嗰幅一流全景。遠處望到啲\VUgras{漁艇仔}
\textsuperscript{(63)}、啲\VUgras{帆船} \textsuperscript{(64)}，仲有一隻大\VUgras{郵輪}
\textsuperscript{(65)}，佢個黑色輪廓喺白\VUgras{雲} \textsuperscript{(67)} 上面凸晒出嚟
——即係\VUgras{地平線} \textsuperscript{(66)} 嗰啲白雲。一陣輕風吹起面\VUgras{旗}
\textsuperscript{(70)}，旗就插喺塊\VUgras{招牌} \textsuperscript{(69)} 上面——\VUgras{大堂}
\textsuperscript{(68)} 嗰塊招牌。

%%K t13-tit-3 sub
\VUsaut{3.0mm}
\VUpk{10.2pt}{40}{消遣。}
\VUsaut{3.0mm}

%%K t13-11-1 p
11. --- 個景實在太好睇，我就決定聽日帶埋\VUgras{觀劇鏡} \textsuperscript{(71)}
或者\VUgras{單筒望遠鏡} \textsuperscript{(72)} 上咖啡室個平台。

%%K t13-12-1 p
12. --- 落咗露台，我行入酒店個咖啡室飲杯\VUgras{嘢} \textsuperscript{(73)}，順便打局
\VUgras{桌球} \textsuperscript{(75)}。我冇停低就穿過個咖啡廳——入面好多人喺度玩
\VUgras{西洋棋} \textsuperscript{(78)}、\VUgras{西洋跳棋} \textsuperscript{(79)}、
\VUgras{骨牌} \textsuperscript{(80)}、\VUgras{十五子棋} \textsuperscript{(81)} 或者
\VUgras{啤牌} \textsuperscript{(82)}——直接上咗\VUgras{桌球房} \textsuperscript{(74)}。

%%K t13-13-1 p
13. --- 打完一局，我落返地下，同侍應攞咗\VUgras{信封} \textsuperscript{(84)}、
\VUgras{信紙} \textsuperscript{(83)}、\VUgras{郵票} \textsuperscript{(85)}、\VUgras{筆}
\textsuperscript{(87)} 同\VUgras{墨水} \textsuperscript{(86)} 嚟寫信。

%%K t13-13-2 p
之後我叫住個\VUgras{馬車伕} \textsuperscript{(92)}，佢架\VUgras{馬車} \textsuperscript{(88)}
就停咗喺咖啡室門口。我問佢係計鐘定係計一程，講掂價之後，佢就幫我開
\VUgras{車門} \textsuperscript{(89)}，扶我入去。佢自己爬返上\VUgras{車頭座}
\textsuperscript{(91)}，一\VUgras{鞭} \textsuperscript{(93)} 就催快啲馬，我哋就出發去遊車河，
一路好舒服。

\VUsaut{6.0mm}
\VUfilet{22mm}
